% CHAPTER 3 - SYSTEM ANALYSIS AND DESIGN
\chapter{System Analysis and Design}

\section{System Architecture Overview}
The Context Memory Framework is a single FastAPI application that serves two interfaces
over the same logic: a JSON API under \texttt{/api} and a server-rendered web interface at
\texttt{/}. Behind them, a set of services implements the pipeline. The application keeps
its own data in a \textit{platform database} (SQLite file \texttt{cmf\_platform.db} by
default) and connects to one or more \textit{target databases} (PostgreSQL, MySQL or
SQLite) that hold the company data being queried. A Qdrant vector store, embedded on disk
unless a server URL is configured, indexes memories, schema tables and past questions.
Google Gemini is the only external service; it generates SQL, writes the plain-language
answer and extracts candidate memories.

\begin{figure}[H]
    \centering
    \begin{tikzpicture}[node distance=0.45cm and 0.5cm, every node/.style={font=\small}]
        \node[blk=5.6cm, fill=gray!10] (ui) {Web UI (Jinja2 + htmx)\quad|\quad JSON API (\texttt{/api})};
        \node[blk=6.6cm, below=of ui, fill=yellow!15] (chat) {Chat pipeline \texttt{/api/chat/ask}\\{\footnotesize auth $\rightarrow$ schema $\rightarrow$ context $\rightarrow$ SQL $\rightarrow$ guard}};
        \node[blk=2.6cm, below left=0.6cm and -1.3cm of chat] (ctx) {Context Builder};
        \node[blk=2.6cm, below right=0.6cm and -1.3cm of chat] (guard) {SQL Guard\\(\texttt{plan\_sql})};
        \node[blk=2.6cm, below=of ctx] (ret) {Memory Retriever\\+ Reranker};
        \node[blk=2.6cm, below=of guard] (chg) {Change Service\\+ Access};
        \node[blk=2.6cm, left=of ret] (store) {Memory Store\\+ Extractor};
        \node[blk=2.6cm, right=of chg] (conn) {Database\\Connector};
        \node[sto, below=of ret] (vec) {Qdrant\\(3 collections)};
        \node[sto, below=of store] (plat) {Platform DB\\(SQLite)};
        \node[sto, below=of conn] (tgt) {Target DBs\\PG / MySQL / SQLite};
        \node[blk=2.2cm, right=1.2cm of chat, fill=red!8] (llm) {Gemini LLM};
        \draw[arr] (ui) -- (chat);
        \draw[arr] (chat) -- (ctx);
        \draw[arr] (chat) -- (guard);
        \draw[arr] (ctx) -- (ret);
        \draw[arr] (guard) -- (chg);
        \draw[arr] (ret) -- (vec);
        \draw[arr] (store) -- (plat);
        \draw[arr] (store) -- (vec);
        \draw[arr] (chg) -- (conn);
        \draw[arr] (conn) -- (tgt);
        \draw[arr, <->] (chat) -- (llm);
        \draw[arr] (ret.west) -- (store.east);
    \end{tikzpicture}
    \caption{High-level architecture of the Context Memory Framework.}
    \label{fig:architecture}
\end{figure}

\section{Proposed Methodology}

\subsection{Architecture Layers}
The system is organized into five logical layers:
\begin{enumerate}[label=\arabic*., leftmargin=*, topsep=4pt, itemsep=4pt]
    \item \textbf{Presentation layer:} page routes in \texttt{app/web/} render Jinja2
          templates and htmx fragments; JSON routers in \texttt{app/api/} serve programmatic
          clients and the test suite. Page routes call the same handler functions as the
          API, so authorization and guardrails exist in one place only.
    \item \textbf{Pipeline layer:} the \texttt{ask} handler orchestrates a single question
          from identification to response.
    \item \textbf{Context layer:} memory storage and extraction, hybrid retrieval with
          reranking, schema-table selection, similar-question examples, and the context
          builder that assembles the prompt.
    \item \textbf{Safety layer:} the SQL planner, the access-control service (roles,
          grants, masking) and the change service that previews, stores and executes
          approved changes.
    \item \textbf{Data layer:} SQLAlchemy models of the platform database, the database
          connector for target databases, and the Qdrant vector store.
\end{enumerate}

\subsection{End-to-End Data Flow}
A question sent to \texttt{POST /api/chat/ask} passes through the following steps:
\begin{enumerate}[label=\arabic*., leftmargin=*, topsep=4pt, itemsep=4pt]
    \item \textbf{Identify:} resolve the organization and user and check that they belong
          together; resolve the target database from a saved \texttt{connection\_id} or
          from inline credentials.
    \item \textbf{Working memory:} load or create the \texttt{Conversation} and store the
          user's \texttt{Message}.
    \item \textbf{Schema:} load the cached schema by a SHA-256 fingerprint of
          type/user/host/port/database, or introspect it (columns, foreign-key references,
          sampled categorical values) and cache it.
    \item \textbf{Context build:} retrieve memories, select schema tables, add similar past
          questions, recent turns and the last \texttt{QueryMemory}, and fit everything into
          the token budget. If this step fails the pipeline falls back to a schema-only
          prompt.
    \item \textbf{Generate SQL:} call the LLM with the built prompt.
    \item \textbf{Plan and guard:} classify the statement as a read or a single-table
          change, or refuse it, and validate it against the live schema.
    \item \textbf{Read:} execute in a transaction that is always rolled back, capped at
          \texttt{MAX\_QUERY\_ROWS}, mask sensitive columns for non-admins, and ask the LLM
          for a plain-language answer (falling back to computed statistics).
          \textbf{Write:} check permission, preview the affected rows, and store a pending
          \texttt{ChangeRequest} for admin approval.
    \item \textbf{Persist:} store the assistant message and a \texttt{QueryMemory} with a
          literal-masked SQL template and structured metadata; index the question as an
          example.
    \item \textbf{Extract (background):} ask the LLM for candidate memories from the turn
          and save them as \textit{pending}, superseding a conflicting memory when the user
          stated a correction.
\end{enumerate}

\begin{figure}[H]
    \centering
    \begin{tikzpicture}[node distance=0.32cm, every node/.style={font=\small}]
        \node[blk=4.2cm] (a) {1--2. Authorize user, load conversation};
        \node[blk=4.2cm, below=of a] (b) {3. Load / introspect schema};
        \node[blk=4.2cm, below=of b, fill=yellow!15] (c) {4. Build context (retrieve + budget)};
        \node[blk=4.2cm, below=of c, fill=red!8] (d) {5. LLM generates SQL};
        \node[diamond, draw, aspect=2.2, below=of d, fill=gray!10, inner sep=1pt] (e) {6. \texttt{plan\_sql}};
        \node[blk=3.6cm, below left=0.5cm and 0.2cm of e] (f) {7a. Read: rolled-back\\transaction, cap, mask};
        \node[blk=3.6cm, below right=0.5cm and 0.2cm of e] (g) {7b. Change: permission,\\preview, pending request};
        \node[blk=3.0cm, right=1.4cm of e] (x) {Refuse with reason\\(\texttt{not\_allowed})};
        \node[blk=3.6cm, below=of f, fill=red!8] (h) {LLM answer + insights};
        \node[blk=3.6cm, below=of h] (i) {8. Save message,\\QueryMemory, example};
        \node[blk=3.6cm, below=of i, fill=green!8] (j) {9. Background memory\\extraction (pending)};
        \draw[arr] (a) -- (b); \draw[arr] (b) -- (c); \draw[arr] (c) -- (d); \draw[arr] (d) -- (e);
        \draw[arr] (e) -| node[lbl, above left]{read} (f);
        \draw[arr] (e) -| node[lbl, below left, pos=0.9]{change} (g);
        \draw[arr] (e) -- node[lbl, above]{invalid} (x);
        \draw[arr] (f) -- (h); \draw[arr] (h) -- (i); \draw[arr] (i) -- (j);
    \end{tikzpicture}
    \caption{Processing of one question by \texttt{POST /api/chat/ask}.}
    \label{fig:dataflow}
\end{figure}

\subsection{Modular Components}
Table~\ref{tab:modules} lists the main modules of the code base and their
responsibilities, with line counts measured on the final source tree.

\begin{table}[H]
    \centering
    \small
    \caption{Modular Component Breakdown of the System}
    \begin{tabular}{p{6.5cm} p{0.9cm} p{7.4cm}}
        \toprule
        \textbf{Module} & \textbf{Lines} & \textbf{Responsibility} \\
        \midrule
        \texttt{api/chat.py}                  & 468 & End-to-end question pipeline, conversations. \\[3pt]
        \texttt{services/retrieval.py}        & 475 & Authorized hybrid ranking, duplicates, conflicts. \\[3pt]
        \texttt{services/context\_builder.py} & 480 & Schema selection, examples, prompt assembly, answer prompt. \\[3pt]
        \texttt{services/memory\_store.py}    & 417 & Memory lifecycle, versions, audit log, ingestion. \\[3pt]
        \texttt{services/vector\_store.py}    & 252 & Qdrant collections, upsert, hybrid RRF search. \\[3pt]
        \texttt{services/changes.py}          & 238 & Change preview, request storage, approval, execution. \\[3pt]
        \texttt{services/sql\_guard.py}       & 229 & SQL planning and read/write classification. \\[3pt]
        \texttt{services/connector.py}        & 187 & Target DB engines, introspection, guarded read/write. \\[3pt]
        \texttt{services/extraction.py}       & 179 & LLM memory extraction, override detection. \\[3pt]
        \texttt{services/indexer.py}          & 179 & Vector index sync and startup rebuild. \\[3pt]
        \texttt{services/embedder.py}         & 137 & Sentence-transformer embeddings (lazy). \\[3pt]
        \texttt{services/reranker.py}         &  78 & Cross-encoder reranking (lazy). \\[3pt]
        \texttt{services/text.py}             &  78 & Tokenizer, follow-up detection, sparse vectors. \\[3pt]
        \texttt{services/access.py}           &  76 & Roles, grants, sensitive-column masking. \\[3pt]
        \texttt{services/llm.py}              &  57 & Gemini client for SQL, answers, extraction. \\[3pt]
        \texttt{web/pages.py}                 & 621 & HTML page routes and htmx fragments. \\
        \bottomrule
    \end{tabular}
    \label{tab:modules}
\end{table}

\section{Data Model}
The platform database is defined with SQLAlchemy 2 declarative models and created at
start-up by \texttt{create\_all} plus a small list of column additions (the project does not
use a migration tool). It contains twelve tables, summarised in
Table~\ref{tab:platform_tables}.

\begin{table}[H]
    \centering
    \small
    \caption{Tables of the Platform Database}
    \begin{tabular}{p{4.6cm} p{10.6cm}}
        \toprule
        \textbf{Table} & \textbf{Purpose and key fields} \\
        \midrule
        \texttt{organizations}        & Tenants; every other record is scoped to one. \\[3pt]
        \texttt{users}                & Members with role \texttt{owner}, \texttt{admin}, \texttt{analyst} or \texttt{viewer}. \\[3pt]
        \texttt{database\_connections}& Saved target databases (type, host, port, name, user, encoded password). \\[3pt]
        \texttt{database\_schemas}    & Cached schema JSON keyed by connection fingerprint. \\[3pt]
        \texttt{conversations}        & Chat sessions with an active-task summary. \\[3pt]
        \texttt{messages}             & User and assistant turns, generated SQL, result metadata, context trace. \\[3pt]
        \texttt{memories}             & Long-term memories (type, content, subject, scope, visibility, importance, confidence, status, validity, provenance, embedding). \\[3pt]
        \texttt{memory\_versions}     & Numbered content history of each memory. \\[3pt]
        \texttt{memory\_audit\_logs}  & Every action on a memory and who performed it. \\[3pt]
        \texttt{query\_memories}      & Structured record of each answered question for follow-ups. \\[3pt]
        \texttt{access\_grants}       & Permission for a member to request a change action on a table. \\[3pt]
        \texttt{change\_requests}     & Proposed data changes with preview, review and outcome. \\
        \bottomrule
    \end{tabular}
    \label{tab:platform_tables}
\end{table}

\begin{figure}[H]
    \centering
    \begin{tikzpicture}[every node/.style={font=\small}]
        \tikzset{er/.style={blk=2.9cm, minimum height=0.8cm, font=\footnotesize}}
        \node[er] (org)  at (0,0)       {organizations};
        \node[er] (usr)  at (3.6,0)     {users};
        \node[er] (dbc)  at (7.6,0)     {database\_connections};
        \node[er] (ds)   at (11.6,0)    {database\_schemas};
        \node[er] (mem)  at (0,-1.8)    {memories};
        \node[er] (conv) at (3.6,-1.8)  {conversations};
        \node[er] (cr)   at (7.6,-1.8)  {change\_requests};
        \node[er] (ag)   at (11.6,-1.8) {access\_grants};
        \node[er] (ver)  at (0,-3.4)    {memory\_versions};
        \node[er] (msg)  at (3.6,-3.4)  {messages};
        \node[er] (qm)   at (7.6,-3.4)  {query\_memories};
        \node[er] (aud)  at (0,-5.0)    {memory\_audit\_logs};
        \draw[arr] (usr) -- (org);
        \draw[arr] (dbc.north) -- ++(0,0.45) -| (org.north);
        \draw[arr, dashed] (ds) -- node[lbl, below]{hash} (dbc);
        \draw[arr] (mem) -- (usr);
        \draw[arr] (conv) -- (usr);
        \draw[arr] (cr) -- (dbc);
        \draw[arr] (ag) -- (dbc);
        \draw[arr] (cr) -- (conv);
        \draw[arr] (ver) -- (mem);
        \draw[arr] (aud.west) -- ++(-0.4,0) |- (mem.west);
        \draw[arr] (msg) -- (conv);
        \draw[arr] (qm) -- (conv);
    \end{tikzpicture}
    \caption{Simplified entity relationships of the platform database. Arrows point to the
    referenced table; links to \texttt{organizations} and \texttt{users} are drawn only once
    for clarity, and the schema cache is linked by fingerprint, not a foreign key.}
    \label{fig:er}
\end{figure}

\section{Memory Model and Lifecycle}

\subsection{Memory Categories}
CMF uses the four memory categories introduced in Chapter~2. Three are long-term and stored
in the \texttt{memories} table; working memory is formed by the conversation and the
structured query record (Table~\ref{tab:memory_types}).

\begin{table}[H]
    \centering
    \small
    \caption{Memory Categories in CMF}
    \begin{tabular}{p{2.0cm} p{6.0cm} p{7.2cm}}
        \toprule
        \textbf{Category} & \textbf{Stored as} & \textbf{Example} \\
        \midrule
        Semantic   & \texttt{memory\_type="semantic"}   & ``Revenue excludes cancelled orders'' \\[3pt]
        Procedural & \texttt{memory\_type="procedural"} & ``Display monetary results in INR'' \\[3pt]
        Episodic   & \texttt{memory\_type="episodic"}   & ``User explored top 10 vendors by outstanding payment'' \\[3pt]
        Working    & \texttt{Conversation}, \texttt{Message}, \texttt{QueryMemory} & Previous question, SQL template, filters, sort, limit \\
        \bottomrule
    \end{tabular}
    \label{tab:memory_types}
\end{table}

Every long-term memory has a \textit{scope} (\texttt{personal}, \texttt{connection} or
\texttt{organization}), a \textit{visibility} (\texttt{personal} or \texttt{shared}), an
importance from 1 to 10, a confidence between 0 and 1, an optional validity interval
(\texttt{valid\_from}, \texttt{valid\_to}), the schema fingerprint and the source
conversation and message. Creating or sharing an organization-wide or shared memory
requires a contributor role (\texttt{owner}, \texttt{admin} or \texttt{analyst}).

\subsection{Lifecycle and Governance}
Figure~\ref{fig:lifecycle} shows the lifecycle. Memories extracted by the LLM always start
as \textit{pending}; only \textit{approved} memories reach the prompt. A memory created
manually through the Knowledge page can be approved immediately. When the user states a
correction, the old memory becomes \textit{superseded} (with \texttt{valid\_to} set) and a
new memory that points to it through \texttt{supersedes\_id} is created. Superseded
memories are retrieved only when the question asks about history, detected by phrases such
as ``previously'', ``historically'', ``used to'' or ``old definition''. Deleting a memory
marks it \textit{deleted} and clears its embedding. Every content edit writes a new
\texttt{MemoryVersion} and every action writes a \texttt{MemoryAuditLog} entry.

\begin{figure}[H]
    \centering
    \begin{tikzpicture}[node distance=0.9cm and 1.9cm, every node/.style={font=\small}]
        \node[st] (p) {pending};
        \node[st, right=of p, fill=blue!10] (a) {approved};
        \node[st, below=of p, fill=red!8] (r) {rejected};
        \node[st, right=of a] (s) {superseded};
        \node[st, below=of a] (e) {expired};
        \node[st, below=of s, fill=gray!15] (d) {deleted};
        \node[left=0.9cm of p, font=\small, align=center] (src) {extracted\\or created};
        \draw[arr] (src) -- (p);
        \draw[arr] (p) -- node[lbl, above]{approve} (a);
        \draw[arr] (p) -- node[lbl, left]{reject} (r);
        \draw[arr] (a) -- node[lbl, above]{correction} (s);
        \draw[arr] (a) -- node[lbl, left]{expire} (e);
        \draw[arr] (a) -- node[lbl, below, sloped]{delete} (d);
        \draw[arr] (src.north) to[bend left=30] node[lbl, above]{auto-approve (manual)} (a.north);
    \end{tikzpicture}
    \caption{Lifecycle states of a long-term memory. Only \textit{approved} memories enter
    prompts; \textit{superseded} ones only for history questions.}
    \label{fig:lifecycle}
\end{figure}

\subsection{One-Time Overrides}
Requests that apply to a single query must not become rules. Before extraction, the latest
user message is checked against regular expressions for phrases such as ``for this report'',
``just this once'', ``only for this'', ``this time only'' and ``temporarily''. If one matches
and the message does not also signal a rule change (``from now on'', ``permanently'',
``definition has changed'', ``update the rule'', and similar), extraction is skipped. The
override still applies to the current query because it is part of the question, which
ranks third in the prompt's priority order.

\section{Retrieval and Scoring Design}

\subsection{Authorization Before Ranking}
Authorization filters are applied before any scoring, both in SQL and in the Qdrant payload
filter: same organization; memory owned by the requester or \texttt{shared}; and either the
same connection or not connection-scoped. Status and validity checks follow, and each
excluded memory is recorded with a reason.

\subsection{Candidate Generation}
Candidates come from a Qdrant hybrid query that fuses a dense search (MiniLM embedding) and
a sparse search (hashed term vectors with IDF weighting) by reciprocal rank fusion, limited
to $\max(\texttt{RAG\_CANDIDATES}, k)$ results. If Qdrant is unavailable the retriever scans
all authorized rows in the platform database.

\subsection{First-Stage Score}
Each candidate memory $m$ for query $q$ receives a weighted score
\begin{equation}
    S(q,m) = w_V V + w_K K + w_R R + w_I I + w_T T
\end{equation}
with default weights $w_V=0.45$, $w_K=0.20$, $w_R=0.15$, $w_I=0.10$, $w_T=0.10$
(Figure~\ref{fig:weights}). The components are:
\begin{align}
    V &= \cos(\mathbf{e}_q, \mathbf{e}_m) \;\text{(clipped to }[0,1]\text{)} \\
    K &= \frac{|\mathcal{T}(q) \cap \mathcal{T}(m)|}{|\mathcal{T}(q)|} \\
    R &= 0.5^{\,a/30}, \quad a = \text{age of } m \text{ in days} \\
    I &= \frac{\text{importance}(m) - 1}{9}
\end{align}
where $\mathcal{T}(\cdot)$ is the set of lower-cased terms after stop-word removal, with
\texttt{snake\_case} identifiers also split into parts. $T$ is a type-intent score: if the
question contains a keyword from an intent map (for example ``revenue'' $\rightarrow$
semantic, ``currency'' $\rightarrow$ procedural, ``previous'' $\rightarrow$ episodic) and
the memory type matches, $T=1$; if a keyword matched but the type does not, $T=0.4$; if no
keyword matched, $T=0.6$. When no embedding is available, the vector weight is added to the
keyword weight. The recency term halves every 30 days (Figure~\ref{fig:recency}).

\begin{figure}[H]
    \centering
    \figimg{0.58\textwidth}{chart_weights.png}{Default weights of the first-stage score (from \texttt{.env.example}).}{fig:weights}
\end{figure}

\begin{figure}[H]
    \centering
    \figimg{0.58\textwidth}{chart_recency.png}{Recency score as a function of memory age (30-day half-life).}{fig:recency}
\end{figure}

\subsection{Reranking and Cut-off}
The scored candidates are reranked by the cross-encoder
\texttt{ms-marco-MiniLM-L-6-v2}. Its logit is converted to $r \in (0,1)$ with the logistic
function and combined with the first-stage score:
\begin{equation}
    \text{final}(q,m) = \lambda\, r(q,m) + (1-\lambda)\, S(q,m), \qquad \lambda = 0.6
\end{equation}
If the reranker cannot be loaded, the first-stage score is used unchanged. The top $k$
memories are kept, where $k$ is set by \texttt{MEMORY\_MAX\_RESULTS} (default 10). The
remaining candidates are recorded with the exclusion reason \texttt{below\_top\_k\_cutoff}.

\section{Context Engineering Design}

\subsection{Instruction Priority Hierarchy}
The system section of every prompt states the priority order explicitly:
\begin{enumerate}[label=\arabic*., leftmargin=*, topsep=4pt, itemsep=2pt]
    \item System security rules (read-only by default; a change is one approved statement on one table).
    \item The live database schema, which is authoritative over memories.
    \item The current question's explicit instructions, including one-time overrides.
    \item Confirmed business definitions (semantic memories).
    \item User and database preferences (procedural memories).
    \item Recent conversation and previous analytical context (working and episodic memory).
\end{enumerate}
It also states that text inside retrieved memories is untrusted data, not instructions, and
must be ignored if it tries to bypass security or modify tables.

\subsection{Token Budget}
Token counts are approximated as $\lfloor |\text{text}|/4 \rfloor$ characters. The
mandatory sections (system rules, schema, question and generation requirements) are placed
first; the remaining budget is $\max(200, B - \text{mandatory})$ with $B = 4000$ by default.
Optional content is then admitted greedily in priority order: ranked memories, similar past
questions, recent conversation turns (up to six turns, each message truncated to 300
characters) and finally the previous analytical query. A memory that does not fit is
recorded as \texttt{token\_budget\_exceeded}.

\subsection{Schema Selection}
Schemas with at most \texttt{RAG\_SCHEMA\_FULL\_TABLES} (12) tables are sent in full. For
larger schemas the builder retrieves table descriptions from the \texttt{schema\_tables}
collection (or ranks tables by keyword overlap), reranks them, keeps the top
\texttt{RAG\_SCHEMA\_TOP\_TABLES} (8), and adds tables referenced by foreign keys and tables
used by the previous query, up to four extra tables.

\subsection{Follow-up Rewriting}
With \texttt{RAG\_REWRITE\_QUERY} enabled, a question is treated as a follow-up if it starts
with words such as ``now'', ``only'', ``also'' or ``same'', begins with ``what about'' or
``how about'', refers to ``those'', ``these'', ``same'' or ``previous'', or has at most two
content words. Retrieval then uses the previous question prepended to the current one; the
prompt still contains the user's exact words.

\section{Safety and Access-Control Design}

\subsection{Statement Planning}
Every generated statement is parsed with \texttt{sqlglot} in the target dialect and
classified by \texttt{plan\_sql}:
\begin{itemize}[leftmargin=*, topsep=4pt, itemsep=3pt]
    \item Exactly one statement is allowed.
    \item A \textbf{read} must be a \texttt{SELECT}, \texttt{UNION} or subquery and must not
          contain any write node or a write keyword (\texttt{DROP}, \texttt{DELETE},
          \texttt{UPDATE}, \texttt{INSERT}, \texttt{ALTER}, \texttt{TRUNCATE},
          \texttt{GRANT}, \texttt{REVOKE}, \texttt{ATTACH}, \texttt{PRAGMA}, \texttt{EXEC}).
    \item A \textbf{change} must be one \texttt{UPDATE}, \texttt{INSERT} or \texttt{DELETE}
          on one table, with no nested write, no joins, \texttt{FROM} or \texttt{USING};
          \texttt{UPDATE} and \texttt{DELETE} need a \texttt{WHERE} clause; assigned
          columns must exist.
    \item Every referenced table must exist in the live schema. Anything else, including
          DDL and administrative commands, is refused with a plain-language reason.
\end{itemize}

\subsection{Roles, Grants and Masking}
Roles \texttt{owner} and \texttt{admin} are administrators, \texttt{analyst} is a
contributor and \texttt{viewer} is read-only. A non-admin may request a change only with an
active \texttt{AccessGrant} for that action on that table (or \texttt{*}) and connection.
Columns whose names contain any configured sensitive pattern (password, secret, token,
api\_key, ssn, card\_number, cvv, pin\_code and others) are masked for non-admins in answers
and previews.

\subsection{Department Roles in the Web Interface}
The web interface demonstrates this model with six roles that a visitor chooses on the
landing page (Table~\ref{tab:roles}). Each role maps to a fixed user of the default
organization. On every switch, \texttt{ensure\_user\_type} creates the user if needed, sets
its platform role and synchronizes its grants: grants for the role's actions on all tables
(\texttt{*}) are added and any other active grant of that user is revoked. Change requests
from department roles stay pending and appear on the administrator's Requests page, while an
administrator's own change requests show Apply and Discard buttons directly in the chat. As
noted in Chapter~1, this role choice is a demonstration and not an authentication boundary.

\begin{table}[H]
    \centering
    \small
    \caption{Roles Offered by the Web Interface}
    \begin{tabular}{p{2.4cm} p{2.4cm} p{4.4cm} p{5.2cm}}
        \toprule
        \textbf{Role} & \textbf{Platform role} & \textbf{Change grants (all tables)} & \textbf{Effect} \\
        \midrule
        Admin      & \texttt{owner}   & (all, as admin)                   & Reads unmasked; applies changes after preview \\[3pt]
        Finance    & \texttt{analyst} & \texttt{update}                   & Requests edits for admin approval \\[3pt]
        Sales      & \texttt{analyst} & \texttt{insert}, \texttt{update}  & Requests additions and edits \\[3pt]
        Operations & \texttt{analyst} & \texttt{insert}, \texttt{update}, \texttt{delete} & Requests any change \\[3pt]
        HR         & \texttt{viewer}  & none                              & Read-only \\[3pt]
        Employee   & \texttt{viewer}  & none                              & Read-only \\
        \bottomrule
    \end{tabular}
    \label{tab:roles}
\end{table}

\subsection{Change Approval}
Figure~\ref{fig:change_flow} shows the change workflow. The approval step re-plans the
stored statement and recounts the matching rows; if the count differs from the preview,
nothing runs. Otherwise the statement executes in a transaction that must affect exactly the
previewed number of rows, or it is rolled back.

\begin{figure}[H]
    \centering
    \begin{tikzpicture}[every node/.style={font=\small}]
        \node[blk=3.0cm] (a) at (0,0)    {Change SQL\\from LLM};
        \node[blk=3.0cm] (b) at (4.6,0)  {Permission:\\admin or grant};
        \node[blk=3.0cm] (c) at (9.2,0)  {Preview: count\\+ first 5 rows};
        \node[blk=3.0cm, fill=yellow!15] (d) at (9.2,-2)  {ChangeRequest\\\textit{pending}};
        \node[blk=3.0cm] (e) at (4.6,-2) {Admin approves:\\re-plan + recount};
        \node[blk=3.0cm, fill=green!8] (f) at (0,-2) {Execute; must touch\\exactly $N$ rows};
        \node[blk=3.6cm, fill=red!8] (g) at (4.6,-4) {failed or rejected:\\nothing changed};
        \draw[arr] (a) -- (b); \draw[arr] (b) -- (c); \draw[arr] (c) -- (d);
        \draw[arr] (d) -- node[lbl, above]{approve} (e);
        \draw[arr] (e) -- node[lbl, above]{match} (f);
        \draw[arr] (e) -- node[lbl, right]{mismatch} (g);
        \draw[arr] (d.south) |- node[lbl, pos=0.75, below]{reject} (g.east);
        \draw[arr] (f.south) |- node[lbl, pos=0.75, below]{row count differs} (g.west);
    \end{tikzpicture}
    \caption{Change-request workflow with re-validation at approval time.}
    \label{fig:change_flow}
\end{figure}

\section{Configuration Parameters}
All settings are defined in \texttt{app/core/config.py} with Pydantic Settings and can be
overridden in \texttt{.env}. The most important values are listed in
Table~\ref{tab:config}; the complete list is in Appendix~\ref{app:config}.

\begin{table}[H]
    \centering
    \small
    \caption{Key System Configuration Parameters (Defaults)}
    \begin{tabular}{p{6.4cm} p{4.2cm} p{4.6cm}}
        \toprule
        \textbf{Parameter} & \textbf{Default} & \textbf{Role} \\
        \midrule
        \texttt{MODEL\_NAME}              & gemini-2.5-flash & SQL generation model \\[1pt]
        \texttt{EMBEDDING\_MODEL}         & all-MiniLM-L6-v2 & 384-d dense embeddings \\[1pt]
        \texttt{RERANKER\_MODEL}          & ms-marco-MiniLM-L-6-v2 & Cross-encoder reranker \\[1pt]
        \texttt{RERANK\_WEIGHT}           & 0.6   & $\lambda$ in final score \\[1pt]
        \texttt{MEMORY\_TOKEN\_BUDGET}    & 4000  & Prompt budget $B$ \\[1pt]
        \texttt{MEMORY\_MAX\_RESULTS}     & 10    & Memories kept after ranking \\[1pt]
        \texttt{RAG\_CANDIDATES}          & 30    & Vector candidates per search \\[1pt]
        \texttt{RAG\_SCHEMA\_FULL\_TABLES}& 12    & Send full schema up to this size \\[1pt]
        \texttt{RAG\_SCHEMA\_TOP\_TABLES} & 8     & Tables kept for large schemas \\[1pt]
        \texttt{RAG\_EXAMPLES}            & 3     & Similar past questions in prompt \\[1pt]
        \texttt{CONVERSATION\_HISTORY\_TURNS} & 6 & Recent turns considered \\[1pt]
        \texttt{MAX\_QUERY\_ROWS}         & 100   & Row cap for reads \\[1pt]
        \texttt{WRITE\_MAX\_ROWS}         & 50    & Row cap for one change \\[1pt]
        \texttt{ALLOW\_DATA\_CHANGES}     & True  & Disable all changes when False \\
        \bottomrule
    \end{tabular}
    \label{tab:config}
\end{table}

\section{System Requirements}

\subsection{Hardware Requirements}
The system has no special hardware requirement. Embedding and reranking models are small
MiniLM models that run on a CPU; a GPU is optional. The development and evaluation machine
used in this project was a laptop with a 12th Gen Intel Core i5-12450HX processor, 16\,GB of
RAM and an NVIDIA GeForce RTX 3050 (6\,GB) laptop GPU, running Windows 11.

\subsection{Software Requirements}
\begin{itemize}[leftmargin=*, topsep=4pt, itemsep=3pt]
    \item Python 3.10 or later.
    \item FastAPI 0.111, Uvicorn 0.30, SQLAlchemy 2.0, Pydantic 2.7 and pydantic-settings.
    \item sqlglot 25.2, qdrant-client 1.12, sentence-transformers 3.0, NumPy and pandas.
    \item google-generativeai 0.7 and a Gemini API key (\texttt{LLM\_API\_KEY}).
    \item Jinja2 and python-multipart for the web interface; htmx 2.0.3 loaded in the browser.
    \item pytest and httpx for testing; psycopg2 for PostgreSQL targets (a Docker Compose file
          for a local PostgreSQL 15 target database is optional).
\end{itemize}
